\section{The complete quadratic calibration}\label{sec:quadratic}

\subsection{Exact energy and both boundary payments}
Take $H_a(z)=z^2-a^2$ on $|z|<R$, with $0<a<R$. In
Lemma~\ref{lem:cylinder}, $m=C=0$ and
\begin{equation}\label{eq:Uround}
 U_a(\theta)=\arsinh\sqrt{(R/a)^2-\cos^2\theta},\qquad
 U_0:=U_a(0)=\arcosh(R/a).
\end{equation}
Here $\omega=\dd u$ and the full energy is the cylinder area.
\begin{proposition}[Exact quadratic calculation]\label{prop:quadratic}
One has
\begin{align}
 \E_a&=2\int_0^{2\pi}U_a(\theta)\dd\theta,
 \label{eq:Equad}\\
 4\pi\arcosh(R/a)&\le\E_a\le4\pi\arsinh(R/a),
 \label{eq:quadbounds}\\
 \E_a&=4\pi\log(2R/a)+O(a^2/R^2)\quad(a/R\longrightarrow0).
 \label{eq:quadasymp}
\end{align}
With $q=2\pi\ii$ and $p=-2U_0$ oriented plus to minus, the real
cut and physical contributions are
\begin{equation}\label{eq:quadpayments}
 E_{\cut}=2\pi U_0,\qquad
 E_{\phys}=\E_a-2\pi U_0.
\end{equation}
Each contains half of the leading logarithmic divergence.
\end{proposition}
\begin{proof}
The coordinate $u=s+\ii\theta$ has energy form
$(\ii/2)\dd u\wedge\dd\bar u=\dd s\dd\theta$, proving
\eqref{eq:Equad}. The minimum and maximum of $U_a$ are
$\arcosh(R/a)$ and $\arsinh(R/a)$, proving \eqref{eq:quadbounds}.
The identities
$\arcosh x=\log(x+\sqrt{x^2-1})$ and
$\arsinh x=\log(x+\sqrt{x^2+1})$ give the stated asymptotics.

To verify every boundary contribution directly, cut the cylinder at
$\theta=0$ and use $F=u$ on $0\le\theta\le2\pi$.
The bottom runs from $-U_0$ to $U_0$, the right side upward,
the top in reverse, and the left side downward. The bottom-to-top jump
is $-2\pi\ii$, so the paired cut integral is $-4\pi\ii U_0$.
Its energy contribution is $2\pi U_0$.
On the right $u=U_a+\ii\theta$ and on the left
$u=-U_a+\ii\theta$, with opposite $\theta$ orientations. Thus
\begin{align*}
 \int_{\partial X}F\bar\omega
 &=2\ii\int_0^{2\pi}(\theta U_a'-U_a)\dd\theta\\
 &=4\pi\ii U_0-4\ii\int_0^{2\pi}U_a\dd\theta.
\end{align*}
Multiplication by $\ii/2$ proves \eqref{eq:quadpayments}. A common
global primitive shift changes neither grouped contribution, since the
integral of $\bar\omega$ over either grouped boundary pair is zero.
\end{proof}

For the heat control $H(\tau,z)=z^2-2\tau$, $a=\sqrt{2\tau}$, so
\begin{equation}\label{eq:quadheat}
 \E_\tau=2\pi\log(1/\tau)+O(1)
 \qquad(\tau\downarrow0).
\end{equation}
The finite core period is not a finite-energy endpoint statement.

\subsection{A scoped limitation of bounded exterior summaries}
\begin{proposition}\label{prop:compression}
Fix $R>0$ and $0<a\le R/2$. The core period, an exterior field floor,
every prescribed finite collection of exterior differential derivative
upper bounds, and separately normalized boundary primitive oscillation
upper bounds can all be chosen uniformly in $a$. Nevertheless $\E_a$
is unbounded as $a\downarrow0$. No finite upper estimate for the energy
can depend only on those fixed common budget inequalities.
\end{proposition}
\begin{proof}
The core period is identically $2\pi\ii$. On $|z|=R$,
$|H_a|\ge3R^2/4$. On a fixed collar, for example
$3R/4\le|z|\le5R/4$, the expression
\[
 \frac1{\sqrt{z^2-a^2}}=\frac1z
       \left(1-\frac{a^2}{z^2}\right)^{-1/2}
\]
is holomorphic with the branch near one and uniformly bounded.
Cauchy's estimates on smaller fixed collars bound every fixed finite
derivative collection. In the cylinder parametrization the separately
rebased boundary primitives have increments
$\pm(U_a(\theta)-U_0)+\ii\theta$; they are uniformly bounded since
\[
 0\le U_a-U_0\le\arsinh(R/a)-\arcosh(R/a).
\]
But \eqref{eq:quadbounds} diverges. The missing relative constant is
$2U_0$, which is unbounded.
\end{proof}

The scope is important. Exact attained norms or complete exact boundary
values can encode $a$ and permit its recovery. The proposition concerns
fixed common upper-budget inequalities, not every possible use of exterior
data. It also shows why an absolute field error set admitting arbitrarily
small positive $a$ admits arbitrarily large finite reciprocal energy,
even if each admitted positive-$a$ fiber is regular.

\subsection{A lower bound for general pair charts}
\begin{lemma}\label{lem:screenlower}
Suppose $H(z)=(z^2-a^2)V(z)$ on $|z|\le R$, $0<2a<R$, and
$|V|\le M$. Then
\begin{equation}\label{eq:lowerlog}
 \E\ge\frac{16\pi}{5M}\log\frac{R}{2a}.
\end{equation}
Consequently, an independently obtained upper $\E\le B$ implies
$a\ge(R/2)\exp[-5MB/(16\pi)]$.
\end{lemma}
\begin{proof}
On $2a\le|z|\le R$, $|z^2-a^2|\le(5/4)|z|^2$.
Integrating $2\cdot4/(5M|z|^2)$ over this annulus proves
\eqref{eq:lowerlog}; rearranging proves the consequence.
\end{proof}
An endpoint conclusion requires a fixed positive $R$ and uniform $M,B$.
The upper estimates later in this manuscript do not supply uniform $B$.

\subsection{Capacity as supporting analysis}
On the round annulus $r<|z|<R$, the function
$\varphi(z)=\log(|z|/r)/\log(R/r)$ has boundary values zero and one.
Direct polar integration gives, with no factor $1/2$,
\[
 \operatorname{cap}A(r,R)=\int|\nabla\varphi|^2\dd A
        =\frac{2\pi}{\log(R/r)}.
\]
For a closed real one-form of core period $p$, circlewise
Cauchy--Schwarz and radial integration give
\[
 \int|\alpha|^2\dd A\ge\frac{p^2}{2\pi}\log(R/r).
\]
Indeed its tangential integral is $p$ on each circle, so its energy on
that circle is at least $p^2/(2\pi s)$. These formulas explain the
degeneration scale but do not give an upper modulus bound. In particular,
a trial potential gives an upper capacity bound and hence a lower bound
for the fixed-period core energy, not the upper estimate sought here.
